\section{字体属性}
\subsection{Family}

有成千上万种字体, 但是基本可以分为三个家族(family), 即

\begin{table}[h]
    \centering
    \begin{tabular}{@{}ll@{}}
    \toprule
    \multicolumn{1}{c}{\textbf{family}} & \multicolumn{1}{c}{\textbf{TeX等价}} \\ \midrule
    有衬线字体(serif)                        & \textbackslash{}rmfamily           \\
    无衬线字体(sans-serif)                   & \textbackslash{}sffamily           \\
    等宽字体(monospaced)~~~~~~           & \textbackslash{}ttfamily           \\ \bottomrule
    \end{tabular}
\end{table}

\subsection{Shape}

每个字体都支持Shape属性和series属性(但是不是全支持, 例如xeCJK就是通过AutoFakeSlant与AutoFakeBold来分别解决确实属性的)。Shape一般有三种：

\begin{table}[h]
    \centering
    \begin{tabular}{@{}ll@{}}
    \toprule
    \multicolumn{1}{c}{\textbf{Shape}} & \multicolumn{1}{c}{\textbf{TeX等价}} \\ \midrule
    正体(upright)                       & \textbackslash{}upshape          \\
    斜体(italic)                  & \textbackslash{}itshape           \\
    倾斜体(slanted)          & \textbackslash{}slshape        \\
    小型大写(Small Capitals)~~~~~~           & \textbackslash{}scshape      \\ \bottomrule
    \end{tabular}
\end{table}

其中\verb|\scshape|理论上对中文无效, 因为中文没有大小写！英文可以用\verb|\lowercase| 或 \verb|\uppercase| 将其转化为小写/大写。

\subsection{Series}

字体的Series属性其实是定义了字体的粗细程度, 它一般分为如下三种：

\begin{table}[h]
    \centering
    \begin{tabular}{@{}ll@{}}
    \toprule
    \multicolumn{1}{c}{\textbf{Series}} & \multicolumn{1}{c}{\textbf{TeX等价}} \\ \midrule
    粗体(bold)                        & \textbackslash{}bfseries          \\
    正常(middle)~~~~~~                  & \textbackslash{}mdseries          \\
    细体(light)         & \textbackslash{}lfseries          \\ \bottomrule
    \end{tabular}
\end{table}

\section{字号}

\subsection{字号命令}

TeX里面设置字体大小的开关从小到大如下：

\begin{multicols}{2}
    \begin{itemize}
      \item \verb|\tiny|
      \item \verb|\scriptsize|
      \item \verb|\footnotesize|
      \item \verb|\small|
      \item \verb|\normalsize|
      \item \verb|\large|
      \item \verb|\Large|
      \item \verb|\Large|
      \item \verb|\huge|
      \item \verb|\Huge|
    \end{itemize}
\end{multicols}

\subsection{中文字号}

使用\verb|\zihao{}|调整字号大小时，西文字号会始终和中文字号保持一致。

\begin{multicols}{4}
    \begin{itemize}
      \item[0] 初号
      \item[-0] 小初
      \item[1] 一号
      \item[-1] 小一
      \item[2] 二号
      \item[-2] 小二
      \item[3] 三号
      \item[-3] 小三
      \item[4] 四号
      \item[-4] 小四
      \item[5] 五号
      \item[-5] 小五
      \item[6] 六号
      \item[-6] 小六
      \item[7] 七号
      \item[8] 八号
    \end{itemize}
\end{multicols}

\section{特殊符号}
\subsection{保留字符}

LaTeX中的命令通常是由一个反斜杠加上命令名称，再加上花括号内的参数构成的。

LaTeX中有许多字符都有特殊的意义，LaTeX中的保留字符有 \verb|#,$,%,&,_,{ }, \|，这些在正文中都不能直接呈现。

反斜杠用\verb|\textbackslash|表示，其他的符号在前面加 \textbackslash \ 便可以表示了。

\subsection{波浪线}

波浪线\~{}也不能在LaTeX中直接输出，因为它的作用是禁止在该处断行，可以尝试：

\begin{lstlisting}
    a $\sim$ b
    a \~ b
    a \~{} b
    a \textasciitilde b
\end{lstlisting}

    a $\sim$ b \qquad 
    a \~ b \qquad 
    a \~{} b \qquad 
    a \textasciitilde b

其中\verb|\sim|为数学公式中的命令故需要加\$\$不然会报错，命令\verb|\~|会自动读取后面的一个字符作为参数，故第二个例子在b上加波浪线，而\verb|\~{}|命令设置参数为空。

\subsection{角度和摄氏度}

角度符号或温度符号（摄氏度）需要借助数学模式。

角度可以使用\verb|$^{\circ}$|，摄氏度可以使用\verb|$^{\circ}\mathrm{C}$|。

还可以添加定义
\begin{lstlisting}
    \def\degree{${}^{\circ}$}
\end{lstlisting}

这样可以直接输入45\verb|\degree|即可。

另外，使用\verb|\usepackage{gensymb}|可以直接使用\verb|\degree|。


\subsection{输出命令}

\verb|verbatim|环境可原样输出其中的文本，忽视TeX命令。

\begin{lstlisting}
    \begin{verbatim}
    \usepackage{times}
    \usepackage{latexsym}
    \end{verbatim}
\end{lstlisting}

输出如下：
\begin{framed}
\begin{verbatim}
    \usepackage{times}
    \usepackage{latexsym}
\end{verbatim}
\end{framed}

对简短的抄录，可使用\verb|\verb||文字|和\verb|\verb*||文字|。

